(1) In the normal form, $v\beta + vw_0\delta \in v\,\Cl(\mathbb{V}')$, and two elements $w_0 + \alpha$, $w_0 + \alpha'$ are equivalent iff $\alpha - \alpha' \in v\,\Cl(\mathbb{V}') \cap \Cl(U') = 0$: by the unique-expansion lemma (Lemma \claimref{transport.splitting-along-hyperbolic-plane}), $v\,\Cl(\mathbb{V}') = v\,\Cl(U') \oplus vw_0\,\Cl(U')$ occupies exactly the $v$- and $vw_0$-slots, which meet the $\Cl(U')$-slot in $0$. (2) Preservation of $S(v)$: since $R_u$ fixes $v$, $v\,(R\sigma R^{-1})\,v = R\,(v\sigma v)\,R^{-1} = RvR^{-1} = v$. Preservation of the ideal: $R\,(vX)\,R^{-1} = (RvR^{-1})(RXR^{-1}) = v\,(RXR^{-1}) \in v\Cl(\mathbb{V}')$. The translation formula is Step 3 of the proof of Theorem \claimref{transport.no-go-level-elements}: conjugating the normal form by $R_u$, every correction retains a left factor of $v$ except the contribution of $R_u w_0 R_u^{-1} = w_0 - u - Q(u)v$, whose $\Cl(U')$-component is $-u$. (3) The orbit map $u \mapsto R_u \cdot [\sigma_0] = [\sigma_0] - u$ is a bijection of $U$ onto the orbit (injectivity by (1)); the cocycle at base $[\sigma_0]$ is $c(R_u) = [\sigma_0] - R_u[\sigma_0] = u = \theta(R_u)$, independent of the base point since the coefficient action is trivial, and $c(R_uR_{u'}) = u + u' = c(R_u) + c(R_{u'})$ is the homomorphism property. The situation is the standard one: an affine action whose linear part is trivial has a fixed point exactly when its translation cocycle is a coboundary, and here the cocycle is tautological, the radical translating the candidates by exactly itself.
